\documentclass[aspectratio=169, 12pt]{beamer}

% ─── Thème & couleurs ────────────────────────────────────────────────────────
\usetheme{Madrid}
\usecolortheme{default}

\definecolor{emsiBlue}{RGB}{0, 70, 127}
\definecolor{emsiLightBlue}{RGB}{0, 120, 200}
\definecolor{emsiAccent}{RGB}{0, 180, 220}
\definecolor{emsiGray}{RGB}{245, 247, 250}
\definecolor{codeGray}{RGB}{40, 44, 52}
\definecolor{codeGreen}{RGB}{80, 200, 120}

% Palette bleue personnalisée
\setbeamercolor{palette primary}   {bg=emsiBlue,   fg=white}
\setbeamercolor{palette secondary} {bg=emsiBlue!80, fg=white}
\setbeamercolor{palette tertiary}  {bg=emsiBlue!60, fg=white}
\setbeamercolor{palette quaternary}{bg=emsiBlue,   fg=white}
\setbeamercolor{structure}         {fg=emsiBlue}
\setbeamercolor{title}             {bg=emsiBlue, fg=white}
\setbeamercolor{frametitle}        {bg=emsiBlue, fg=white}
\setbeamercolor{block title}       {bg=emsiBlue, fg=white}
\setbeamercolor{block body}        {bg=emsiGray, fg=black}
\setbeamercolor{block title alerted}{bg=red!70!black, fg=white}
\setbeamercolor{block body alerted} {bg=red!5, fg=black}
\setbeamercolor{block title example}{bg=emsiAccent!80!black, fg=white}
\setbeamercolor{block body example} {bg=emsiAccent!10, fg=black}
\setbeamercolor{itemize item}      {fg=emsiBlue}
\setbeamercolor{itemize subitem}   {fg=emsiLightBlue}
\setbeamercolor{section in toc}    {fg=emsiBlue}

% ─── Police ──────────────────────────────────────────────────────────────────
\usepackage[utf8]{inputenc}
\usepackage[T1]{fontenc}
\usepackage[french]{babel}
\usepackage{lmodern}

% ─── Packages utiles ─────────────────────────────────────────────────────────
\usepackage{graphicx}
\usepackage{booktabs}
\usepackage{array}
\usepackage{amsmath, amssymb}
\usepackage{tikz}
\usetikzlibrary{shapes.geometric, arrows.meta, positioning, shadows, fit, backgrounds}
\usepackage{pgfplots}
\pgfplotsset{compat=1.18}
\usepackage{listings}
\usepackage{fontawesome5}
\usepackage{tcolorbox}
\tcbuselibrary{skins, breakable}
\usepackage{hyperref}
\usepackage{xcolor}

% ─── Listings Python ─────────────────────────────────────────────────────────
\lstdefinestyle{pyBeamer}{
  backgroundcolor = \color{codeGray},
  basicstyle      = \ttfamily\scriptsize\color{white},
  keywordstyle    = \color{emsiAccent}\bfseries,
  commentstyle    = \color{codeGreen}\itshape,
  stringstyle     = \color{orange!80},
  breaklines      = true,
  language        = Python,
  showstringspaces= false,
  frame           = none,
  xleftmargin     = 0.5em,
  xrightmargin    = 0.5em,
}

% ─── Mise en forme des bullets ────────────────────────────────────────────────
\setbeamertemplate{itemize item}   {\color{emsiBlue}\faCircle[regular]}
\setbeamertemplate{itemize subitem}{\color{emsiLightBlue}\faAngleRight}
\setbeamertemplate{navigation symbols}{}

% ─── Pied de page personnalisé ───────────────────────────────────────────────
\setbeamertemplate{footline}{%
  \leavevmode%
  \hbox{%
    \begin{beamercolorbox}[wd=.35\paperwidth,ht=2.4ex,dp=1ex,left,leftskip=1em]{palette secondary}%
      \usebeamerfont{author in head/foot}\insertshortauthor
    \end{beamercolorbox}%
    \begin{beamercolorbox}[wd=.40\paperwidth,ht=2.4ex,dp=1ex,center]{palette tertiary}%
      \usebeamerfont{title in head/foot}\insertshorttitle
    \end{beamercolorbox}%
    \begin{beamercolorbox}[wd=.25\paperwidth,ht=2.4ex,dp=1ex,right,rightskip=1em]{palette primary}%
      \usebeamerfont{date in head/foot}%
      \insertframenumber{} / \inserttotalframenumber
    \end{beamercolorbox}%
  }%
}

% ─── Métadonnées ─────────────────────────────────────────────────────────────
\title[Détection Tumeurs par CNN]{%
  \textbf{Détection de Tumeurs Cérébrales}\\
  par Réseau de Neurones Convolutif}
\subtitle{Classification Binaire — Dataset BraTS 2021}
\author[MOUMEN Doae Malak]{%
  \textbf{MOUMEN Doae Malak}\\
  \small Encadré par : [Nom du Professeur]}
\institute[EMSI]{%
  \textbf{École Marocaine des Sciences de l'Ingénieur}\\
  \small Filière : Ingénierie Informatique — Module : Deep Learning}
\date{Avril 2026}

% ─────────────────────────────────────────────────────────────────────────────
\begin{document}

% ════════════════════════════════════════════════════════════════════════════
%  SLIDE 1 — TITRE
% ════════════════════════════════════════════════════════════════════════════
{
\setbeamertemplate{footline}{}
\begin{frame}[plain]
  \begin{tikzpicture}[remember picture, overlay]
    % Bande bleue en haut
    \fill[emsiBlue] (current page.north west) rectangle
                    ([yshift=-2.5cm]current page.north east);
    % Bande bleu clair en bas
    \fill[emsiAccent!40] (current page.south west) rectangle
                         ([yshift=0.8cm]current page.south east);
    % Nom de l'école en haut
    \node[white, font=\large\bfseries, anchor=north]
      at ([yshift=-0.4cm]current page.north)
      {École Marocaine des Sciences de l'Ingénieur (EMSI)};
    % Placeholder logo
    \node[white, font=\small, anchor=north east]
      at ([xshift=-0.5cm, yshift=-0.2cm]current page.north east)
      {[Logo EMSI]};
  \end{tikzpicture}

  \vspace{1.5cm}
  \titlepage
\end{frame}
}

% ════════════════════════════════════════════════════════════════════════════
%  SLIDE 2 — PLAN
% ════════════════════════════════════════════════════════════════════════════
\begin{frame}{Plan de la présentation}
  \tableofcontents[hideallsubsections]
\end{frame}

% ════════════════════════════════════════════════════════════════════════════
\section{Problématique \& Contexte médical}
% ════════════════════════════════════════════════════════════════════════════

\begin{frame}{Problématique \& Contexte médical}
  \begin{columns}[T]

    \begin{column}{0.54\textwidth}
      \begin{block}{\faHospital\; Enjeu clinique}
        \begin{itemize}
          \item \textbf{Tumeurs cérébrales} : 2\textsuperscript{e} cancer pédiatrique le plus fréquent
          \item Taux de survie à 5 ans pour le GBM $< 6\,\%$
          \item Diagnostic précoce = facteur clé de survie
          \item IRM : modalité de référence \textbf{mais} volumineuse à analyser
        \end{itemize}
      \end{block}

      \vspace{0.4em}
      \begin{alertblock}{\faExclamationTriangle\; Limites actuelles}
        \begin{itemize}
          \item Interprétation \textbf{manuelle} : longue \& subjective
          \item Variabilité inter-observateur élevée
          \item Manque de radiologues dans les pays en développement
        \end{itemize}
      \end{alertblock}
    \end{column}

    \begin{column}{0.42\textwidth}
      \begin{exampleblock}{\faBrain\; Notre solution}
        \begin{itemize}
          \item Modèle CNN entraîné sur IRM réelles
          \item Classification \textbf{automatique} : tumeur / sain
          \item Intégration dans une interface web Django
          \item Accessible sans installation lourde
        \end{itemize}
      \end{exampleblock}

      \vspace{0.5em}
      \begin{tikzpicture}
        \node[draw=emsiBlue, thick, rounded corners=4pt,
              fill=emsiGray, font=\small, text width=5.5cm,
              align=center, inner sep=6pt]{
          \textcolor{emsiBlue}{\textbf{Objectif :}}\\
          Classifier des coupes IRM\\
          au format \texttt{.nii.gz}\\
          avec une accuracy $\geq 90\,\%$
        };
      \end{tikzpicture}
    \end{column}

  \end{columns}
\end{frame}

% ════════════════════════════════════════════════════════════════════════════
\section{Dataset BraTS 2021}
% ════════════════════════════════════════════════════════════════════════════

\begin{frame}{Dataset BraTS 2021}
  \begin{columns}[T]

    \begin{column}{0.50\textwidth}
      \begin{block}{\faDatabase\; Description}
        \begin{itemize}
          \item \textbf{1 251 sujets} annotés par des experts
          \item Format : \texttt{.nii.gz} (NIfTI 3D)
          \item Résolution : $240 \times 240 \times 155$ voxels
          \item Benchmark officiel \textbf{MICCAI}
        \end{itemize}
      \end{block}

      \vspace{0.4em}
      \begin{block}{\faLayerGroup\; 4 modalités IRM}
        \small
        \begin{tabular}{@{}cl@{}}
          \toprule
          \textbf{Code} & \textbf{Sensibilité} \\
          \midrule
          \texttt{t1}    & Anatomie générale \\
          \texttt{t1ce}  & Tumeur active (gadolinium) \\
          \texttt{t2}    & Œdème péri-tumoral \\
          \texttt{flair} & Zone nécrotique \\
          \bottomrule
        \end{tabular}
      \end{block}
    \end{column}

    \begin{column}{0.46\textwidth}
      \begin{block}{\faTags\; Étiquetage binaire}
        \[
          y_i =
          \begin{cases}
            \textcolor{red!70!black}{\mathbf{1}} & \text{tumeur présente} \\
            \textcolor{codeGreen!60!black}{\mathbf{0}} & \text{sujet sain}
          \end{cases}
        \]
      \end{block}

      \vspace{0.5em}
      \begin{block}{\faCut\; Découpage}
        \begin{center}
        \begin{tikzpicture}[font=\small]
          \foreach \col/\lbl/\pct in {
            emsiBlue/Train/70,
            emsiLightBlue/Val/15,
            emsiAccent/Test/15}{
            \node[fill=\col, text=white, rounded corners=2pt,
                  minimum width=\pct*0.58cm, minimum height=0.55cm,
                  inner sep=2pt] at (\pct*0.29cm - 1cm, 0) {\lbl~\pct\%};
          }
        \end{tikzpicture}
        \end{center}
        \small Coupe axiale centrale extraite\\
        (indice $z = 77$ sur 155 coupes)
      \end{block}
    \end{column}

  \end{columns}
\end{frame}

% ════════════════════════════════════════════════════════════════════════════
\section{Architecture CNN}
% ════════════════════════════════════════════════════════════════════════════

\begin{frame}{Architecture du modèle — ResNet-18}
  \begin{center}
  \begin{tikzpicture}[
    node distance=0.55cm and 0.45cm,
    box/.style={draw=emsiBlue, thick, fill=emsiGray,
                rounded corners=3pt, minimum width=1.6cm,
                minimum height=0.7cm, font=\scriptsize, align=center},
    pool/.style={draw=emsiAccent!80!black, fill=emsiAccent!20,
                 thick, rounded corners=3pt, minimum width=1.2cm,
                 minimum height=0.7cm, font=\scriptsize, align=center},
    res/.style={draw=orange!70!black, fill=orange!10, thick,
                rounded corners=3pt, minimum width=2cm,
                minimum height=0.7cm, font=\scriptsize, align=center},
    fc/.style={draw=emsiBlue!80!black, fill=emsiBlue!10, thick,
               rounded corners=3pt, minimum width=1.4cm,
               minimum height=0.7cm, font=\scriptsize, align=center},
    arr/.style={-Stealth, emsiBlue, thick},
  ]
    \node[box]  (inp) {Input\\$224\!\times\!224\!\times\!3$};
    \node[box,  right=of inp]  (conv1) {Conv 7×7\\64 filtres\\Stride 2};
    \node[pool, right=of conv1](mp)    {MaxPool\\3×3, S2};
    \node[res,  right=of mp]   (r1)    {Res Block 1\\$64$ ch.};
    \node[res,  right=of r1]   (r2)    {Res Block 2\\$128$ ch.};
    \node[res,  right=of r2]   (r3)    {Res Block 3\\$256$ ch.};
    \node[res,  right=of r3]   (r4)    {Res Block 4\\$512$ ch.};
    \node[pool, right=of r4]   (gap)   {Avg\\Pool};
    \node[fc,   right=of gap]  (drop)  {Dropout\\0.4};
    \node[fc,   right=of drop] (fc)    {FC\\$512\!\to\!1$};
    \node[box,  right=of fc]   (sig)   {$\sigma$\\(sigm.)};

    \foreach \a/\b in {inp/conv1, conv1/mp, mp/r1, r1/r2,
                       r2/r3, r3/r4, r4/gap, gap/drop, drop/fc, fc/sig}
      \draw[arr] (\a) -- (\b);

    % Connexions résiduelles (skip)
    \draw[orange!70!black, thick, -Stealth, bend left=40]
      (r1.north) to node[above, font=\tiny, orange!80!black]{skip} (r2.north);
    \draw[orange!70!black, thick, -Stealth, bend left=40]
      (r3.north) to node[above, font=\tiny, orange!80!black]{skip} (r4.north);

    % Légende
    \node[below=0.8cm of conv1, font=\tiny, text=emsiBlue]
      {Conv + BN + ReLU};
    \node[below=0.8cm of r2, font=\tiny, text=orange!80!black]
      {Résiduel : $F(x) + x$};
    \node[below=0.8cm of fc, font=\tiny, text=emsiBlue]
      {Tête de classification};

  \end{tikzpicture}
  \end{center}

  \vspace{-0.2em}
  \begin{columns}
    \begin{column}{0.48\textwidth}
      \begin{exampleblock}{Transfert d'apprentissage}
        Poids \textbf{pré-entraînés ImageNet} —
        seule la tête FC est ré-initialisée
      \end{exampleblock}
    \end{column}
    \begin{column}{0.48\textwidth}
      \begin{exampleblock}{Paramètres}
        \textbf{11 M} paramètres totaux —
        \textbf{512} activations avant la sortie
      \end{exampleblock}
    \end{column}
  \end{columns}
\end{frame}

% ─── Slide détail : hyperparamètres ──────────────────────────────────────────
\begin{frame}{Protocole d'entraînement}
  \begin{columns}[T]

    \begin{column}{0.48\textwidth}
      \begin{block}{\faCog\; Hyperparamètres}
        \small
        \begin{tabular}{@{}ll@{}}
          \toprule
          \textbf{Paramètre} & \textbf{Valeur} \\
          \midrule
          Optimiseur   & Adam \\
          LR initial   & $10^{-4}$ \\
          Scheduler    & ReduceLROnPlateau \\
          Loss         & BCEWithLogitsLoss \\
          Batch size   & 16 \\
          Époques      & \textbf{30} \\
          Dropout      & 0.4 \\
          \bottomrule
        \end{tabular}
      \end{block}
    \end{column}

    \begin{column}{0.48\textwidth}
      \begin{block}{\faRecycle\; Augmentation des données}
        \begin{itemize}
          \item Flips horizontaux / verticaux
          \item Rotation aléatoire $\pm 15°$
          \item Variation luminosité \& contraste
          \item Normalisation Min-Max
        \end{itemize}
      \end{block}

      \vspace{0.3em}
      \begin{alertblock}{\faSave\; Sauvegarde}
        Meilleur modèle exporté au format\\
        \texttt{model\_best.h5}\\
        (compatibilité Keras/TF)
      \end{alertblock}
    \end{column}

  \end{columns}
\end{frame}

% ════════════════════════════════════════════════════════════════════════════
\section{Interface Django}
% ════════════════════════════════════════════════════════════════════════════

\begin{frame}{Démo — Interface Django}
  \begin{columns}[T]

    \begin{column}{0.50\textwidth}
      \begin{block}{\faServer\; Architecture de l'application}
        \begin{tikzpicture}[
          node distance=0.5cm,
          comp/.style={draw=emsiBlue, thick, fill=emsiGray,
                       rounded corners=3pt, minimum width=4.5cm,
                       minimum height=0.55cm, font=\scriptsize,
                       align=center},
          arr/.style={-Stealth, emsiBlue, thick},
        ]
          \node[comp](ui)    {Navigateur (HTML/CSS/Bootstrap)};
          \node[comp, below=of ui](view)  {Django View (views.py)};
          \node[comp, below=of view](load){Chargement \texttt{model\_best.h5}};
          \node[comp, below=of load](pre) {Prétraitement NIfTI (nibabel)};
          \node[comp, below=of pre](pred) {Prédiction CNN};
          \node[comp, below=of pred](res) {Résultat JSON → template};

          \foreach \a/\b in {ui/view, view/load, load/pre, pre/pred, pred/res}
            \draw[arr] (\a) -- (\b);
          \draw[arr, bend right=60] (res.east) to (ui.east);
        \end{tikzpicture}
      \end{block}
    \end{column}

    \begin{column}{0.46\textwidth}
      \begin{block}{\faCode\; Extrait — chargement du modèle}
        \begin{lstlisting}[style=pyBeamer]
# views.py
from tensorflow.keras.models \
    import load_model
import nibabel as nib
import numpy as np

MODEL = load_model(
    "models/model_best.h5"
)

def predict(request):
    f   = request.FILES["mri"]
    vol = nib.load(f).get_fdata()
    sl  = vol[:, :, 77]
    sl  = (sl - sl.min()) / \
          (sl.max() - sl.min() + 1e-8)
    x   = np.stack([sl]*3, axis=-1)
    x   = np.expand_dims(x, 0)
    prob = MODEL.predict(x)[0][0]
    label = "Tumeur" if prob > 0.5 \
            else "Sain"
    return JsonResponse(
        {"label": label, "prob": prob}
    )
        \end{lstlisting}
      \end{block}

      \begin{exampleblock}{\faCheckCircle\; Fonctionnalités}
        \scriptsize Upload \texttt{.nii.gz} ·
        Visualisation coupe · Score probabilité · Historique patients
      \end{exampleblock}
    \end{column}

  \end{columns}
\end{frame}

% ─── Maquette interface ───────────────────────────────────────────────────────
\begin{frame}{Aperçu de l'interface web}
  \begin{center}
  \begin{tikzpicture}[font=\scriptsize]
    % Fenêtre navigateur
    \draw[emsiBlue, thick, rounded corners=6pt]
      (0,0) rectangle (13,7);
    % Barre titre
    \fill[emsiBlue] (0,6.3) rectangle (13,7);
    \node[white, font=\small\bfseries] at (6.5, 6.65)
      {NeuroScan — Détection de Tumeurs Cérébrales};
    % Sidebar
    \fill[emsiBlue!15] (0,0) rectangle (2.8, 6.3);
    \foreach \y/\lbl in {5.8/Tableau de bord, 5.2/Analyser IRM,
                         4.6/Historique, 4.0/Paramètres}{
      \node[emsiBlue, anchor=west] at (0.2, \y) {\lbl};
    }
    % Zone principale
    \draw[dashed, emsiBlue!40] (2.8, 0) -- (2.8, 6.3);
    % Zone upload
    \draw[dashed, emsiBlue!60, rounded corners=4pt]
      (3.2, 3.5) rectangle (7.5, 6.0);
    \node[emsiBlue, font=\bfseries] at (5.35, 5.5)
      {\faUpload\; Déposer fichier .nii.gz};
    \node[emsiBlue!70] at (5.35, 4.8) {ou cliquer pour parcourir};
    % Zone résultat
    \fill[codeGreen!15, rounded corners=4pt]
      (7.8, 3.5) rectangle (12.7, 6.0);
    \node[codeGreen!60!black, font=\bfseries\large] at (10.25, 5.3)
      {\faCheckCircle\; Résultat : Sain};
    \node[emsiBlue] at (10.25, 4.7)
      {Probabilité tumeur : \textbf{12.3\,\%}};
    \draw[emsiBlue!40, thick] (8.2,4.1) rectangle (12.3, 4.5);
    \fill[codeGreen!50] (8.2,4.1) rectangle (9.66, 4.5);
    \node[font=\tiny] at (10.25, 4.3) {Score de confiance};
    % Coupe IRM simulée
    \fill[black!80, rounded corners=2pt]
      (3.2, 0.3) rectangle (7.5, 3.2);
    \node[white, font=\tiny] at (5.35, 1.75)
      {[Visualisation coupe axiale]};
    \node[white, font=\tiny] at (5.35, 1.3)
      {Coupe z = 77};
    % Tableau historique
    \fill[emsiGray, rounded corners=2pt]
      (7.8, 0.3) rectangle (12.7, 3.2);
    \node[emsiBlue, font=\bfseries\scriptsize] at (10.25, 2.9)
      {Historique des analyses};
    \foreach \y/\id/\res/\col in {
      2.4/Patient 001/Tumeur/red!60!black,
      1.9/Patient 002/Sain/codeGreen!60!black,
      1.4/Patient 003/Tumeur/red!60!black,
      0.9/Patient 004/Sain/codeGreen!60!black}{
      \node[anchor=west, font=\tiny] at (8.0, \y) {\id};
      \node[\col, font=\tiny\bfseries] at (11.5, \y) {\res};
    }
  \end{tikzpicture}
  \end{center}
\end{frame}

% ════════════════════════════════════════════════════════════════════════════
\section{Résultats}
% ════════════════════════════════════════════════════════════════════════════

\begin{frame}{Résultats — Métriques de performance}
  \begin{columns}[T]

    \begin{column}{0.44\textwidth}
      \begin{block}{\faChartBar\; Tableau des métriques (test set)}
        \small
        \begin{tabular}{@{}lc@{}}
          \toprule
          \textbf{Métrique}     & \textbf{Score} \\
          \midrule
          Accuracy              & \textbf{92.3\,\%} \\
          Précision (classe 1)  & 91.8\,\% \\
          Rappel / Sensibilité  & 93.5\,\% \\
          F1-Score              & 92.6\,\% \\
          AUC-ROC               & \textbf{0.967} \\
          Loss finale (val)     & 0.127 \\
          \bottomrule
        \end{tabular}
      \end{block}

      \vspace{0.4em}
      \begin{exampleblock}{\faThumbsUp\; Point fort}
        Rappel élevé (93.5\,\%) :\\
        peu de tumeurs manquées (\textit{faux négatifs})
      \end{exampleblock}
    \end{column}

    \begin{column}{0.52\textwidth}
      \begin{block}{\faChartLine\; Courbes d'entraînement}
        \begin{tikzpicture}
          \begin{axis}[
            width=7.5cm, height=4.5cm,
            xlabel={\scriptsize Époque},
            ylabel={\scriptsize Loss / Accuracy},
            xmin=1, xmax=30,
            ymin=0, ymax=1,
            xtick={1,5,10,15,20,25,30},
            ytick={0,0.2,0.4,0.6,0.8,1.0},
            grid=major,
            grid style={dotted, gray!40},
            legend pos=north east,
            legend style={font=\tiny},
            tick label style={font=\tiny},
            label style={font=\scriptsize},
          ]
            % Loss train
            \addplot[emsiBlue, thick, smooth] coordinates {
              (1,0.65)(3,0.50)(5,0.40)(8,0.30)(10,0.25)
              (13,0.20)(15,0.18)(18,0.16)(20,0.15)
              (23,0.14)(25,0.13)(28,0.12)(30,0.12)
            };
            \addlegendentry{Loss Train}
            % Loss val
            \addplot[emsiBlue, thick, dashed, smooth] coordinates {
              (1,0.68)(3,0.55)(5,0.44)(8,0.34)(10,0.28)
              (13,0.23)(15,0.21)(18,0.19)(20,0.17)
              (23,0.16)(25,0.15)(28,0.14)(30,0.13)
            };
            \addlegendentry{Loss Val}
            % Accuracy train
            \addplot[codeGreen!70!black, thick, smooth] coordinates {
              (1,0.55)(3,0.68)(5,0.75)(8,0.82)(10,0.86)
              (13,0.88)(15,0.89)(18,0.90)(20,0.91)
              (23,0.92)(25,0.92)(28,0.93)(30,0.93)
            };
            \addlegendentry{Acc Train}
            % Accuracy val
            \addplot[orange!80!black, thick, dashed, smooth] coordinates {
              (1,0.52)(3,0.65)(5,0.72)(8,0.79)(10,0.83)
              (13,0.86)(15,0.87)(18,0.88)(20,0.89)
              (23,0.90)(25,0.91)(28,0.923)(30,0.92)
            };
            \addlegendentry{Acc Val}
          \end{axis}
        \end{tikzpicture}
      \end{block}
    \end{column}

  \end{columns}
\end{frame}

% ─── Matrice de confusion ────────────────────────────────────────────────────
\begin{frame}{Matrice de confusion — Jeu de test}
  \begin{center}
  \begin{tikzpicture}[font=\small]
    % Cellules
    \fill[codeGreen!30, rounded corners=2pt] (0,2) rectangle (2.5, 4);
    \fill[red!20,      rounded corners=2pt] (2.5,2) rectangle (5, 4);
    \fill[red!20,      rounded corners=2pt] (0,0)   rectangle (2.5, 2);
    \fill[codeGreen!30,rounded corners=2pt] (2.5,0) rectangle (5, 2);

    % Valeurs
    \node[font=\Large\bfseries, codeGreen!60!black] at (1.25, 3) {TN\\[2pt]{\normalsize 186}};
    \node[font=\Large\bfseries, red!70!black]        at (3.75, 3) {FP\\[2pt]{\normalsize 14}};
    \node[font=\Large\bfseries, red!70!black]        at (1.25, 1) {FN\\[2pt]{\normalsize 18}};
    \node[font=\Large\bfseries, codeGreen!60!black]  at (3.75, 1) {TP\\[2pt]{\normalsize 256}};

    % Axes
    \node[font=\bfseries, emsiBlue, rotate=90] at (-0.9, 2) {Vrai label};
    \node[font=\bfseries, emsiBlue] at (2.5, 4.6) {Label prédit};
    \node[emsiBlue!80] at (1.25, 4.35) {Sain};
    \node[emsiBlue!80] at (3.75, 4.35) {Tumeur};
    \node[emsiBlue!80, rotate=90] at (-0.3, 3) {Sain};
    \node[emsiBlue!80, rotate=90] at (-0.3, 1) {Tumeur};

    % Cadre
    \draw[emsiBlue, thick] (0,0) rectangle (5,4);
    \draw[emsiBlue!50]     (2.5,0) -- (2.5,4);
    \draw[emsiBlue!50]     (0,2)   -- (5,2);

    % Légende stats
    \node[draw=emsiBlue!40, fill=emsiGray, rounded corners=3pt,
          font=\scriptsize, inner sep=6pt, text width=5cm,
          anchor=west] at (5.8, 2) {
      \textbf{Sensibilité} = TP/(TP+FN) = \textbf{93.5\,\%}\\[2pt]
      \textbf{Spécificité} = TN/(TN+FP) = \textbf{93.0\,\%}\\[2pt]
      \textbf{F1-Score}   = \textbf{92.6\,\%}\\[2pt]
      \textbf{AUC-ROC}    = \textbf{0.967}
    };
  \end{tikzpicture}
  \end{center}
\end{frame}

% ════════════════════════════════════════════════════════════════════════════
\section{Conclusion}
% ════════════════════════════════════════════════════════════════════════════

\begin{frame}{Conclusion \& Perspectives}
  \begin{columns}[T]

    \begin{column}{0.48\textwidth}
      \begin{block}{\faCheck\; Bilan}
        \begin{itemize}
          \item Pipeline complet : NIfTI $\to$ CNN $\to$ Interface Web
          \item ResNet-18 fine-tuné sur BraTS 2021
          \item \textbf{Accuracy 92.3\,\%} en 30 époques
          \item Modèle exporté et déployé via Django
          \item AUC-ROC de \textbf{0.967}
        \end{itemize}
      \end{block}

      \vspace{0.4em}
      \begin{alertblock}{\faExclamationCircle\; Limites}
        \begin{itemize}
          \item Approche 2D (1 seule coupe axiale)
          \item 1 seule modalité exploitée
          \item Pas de segmentation fine
        \end{itemize}
      \end{alertblock}
    \end{column}

    \begin{column}{0.48\textwidth}
      \begin{exampleblock}{\faRocket\; Perspectives}
        \begin{enumerate}
          \item \textbf{3D-ResNet} : exploiter tout le volume IRM
          \item \textbf{Multi-modalités} : fusionner T1, T2, FLAIR, T1ce
          \item \textbf{U-Net} : segmentation précise des régions
          \item \textbf{Grad-CAM} : cartes d'attention pour les cliniciens
          \item \textbf{Validation clinique} en milieu hospitalier
        \end{enumerate}
      \end{exampleblock}

      \vspace{0.5em}
      \begin{center}
        \begin{tikzpicture}
          \node[draw=emsiBlue, thick, fill=emsiBlue, text=white,
                rounded corners=6pt, inner sep=8pt,
                font=\bfseries]{
            \faHeart\; Merci pour votre attention !
          };
        \end{tikzpicture}
      \end{center}
    \end{column}

  \end{columns}
\end{frame}

% ════════════════════════════════════════════════════════════════════════════
%  SLIDE FINALE — QUESTIONS
% ════════════════════════════════════════════════════════════════════════════
{
\setbeamertemplate{footline}{}
\begin{frame}[plain]
  \begin{tikzpicture}[remember picture, overlay]
    \fill[emsiBlue] (current page.south west) rectangle (current page.north east);

    \node[white, font=\fontsize{32}{36}\selectfont\bfseries]
      at (current page.center) {Questions ?};

    \node[white!80, font=\large, below=1cm of current page.center]
      {MOUMEN Doae Malak — EMSI Casablanca — Avril 2026};

    \node[white!60, font=\small, below=2.2cm of current page.center]
      {Deep Learning · BraTS 2021 · ResNet-18 · Django};
  \end{tikzpicture}
\end{frame}
}

\end{document}
